\documentclass[11pt,a4paper]{article}

\usepackage[utf8]{inputenc}
\usepackage[T1]{fontenc}
\usepackage{lmodern}
\usepackage[margin=0.7in]{geometry}
\usepackage[table]{xcolor}
\usepackage{array}
\usepackage{tabularx}
\usepackage{booktabs}
\usepackage{enumitem}
\usepackage{titlesec}
\usepackage{fancyhdr}
\usepackage{lastpage}
\usepackage{hyperref}
\usepackage[normalem]{ulem}

\definecolor{TextInk}{HTML}{233142}
\definecolor{HeadingBlue}{HTML}{1D3557}
\definecolor{AccentTeal}{HTML}{2A7F92}
\definecolor{AccentGold}{HTML}{C58B4E}
\definecolor{RuleGray}{HTML}{D6E0E5}
\definecolor{TableStripe}{HTML}{F3F8FA}
\definecolor{SoftTint}{HTML}{F7F4EE}
\definecolor{SoftBlue}{HTML}{EEF6F8}
\definecolor{SoftGreen}{HTML}{EEF7F0}
\definecolor{SoftRed}{HTML}{FBF0EF}

\hypersetup{
  colorlinks=true,
  linkcolor=HeadingBlue,
  urlcolor=AccentTeal,
  pdftitle={IB Geography SL Food and Health Essay Planning Workbook - 2026-04-30},
  pdfauthor={IB Geography SL Preparation}
}

\color{TextInk}
\setlength{\parskip}{0.45em}
\setlength{\parindent}{0pt}
\setlength{\tabcolsep}{6pt}
\renewcommand{\arraystretch}{1.25}
\setlist[itemize]{leftmargin=1.3em,itemsep=2pt,topsep=3pt}
\setlist[enumerate]{leftmargin=1.45em,itemsep=3pt,topsep=4pt}

\newcolumntype{Y}{>{\raggedright\arraybackslash}X}
\newcolumntype{L}[1]{>{\raggedright\arraybackslash}p{#1}}

\titleformat{\section}
  {\normalfont\Large\sffamily\bfseries\color{HeadingBlue}}
  {}{0pt}{}
  [\vspace{0.1em}{\color{AccentGold}\titlerule[0.8pt]}]

\titleformat{\subsection}
  {\normalfont\large\sffamily\bfseries\color{AccentTeal}}
  {}{0pt}{}

\pagestyle{fancy}
\setlength{\headheight}{14pt}
\fancyhf{}
\fancyhead[L]{\sffamily\color{AccentTeal}April 30 Workbook}
\fancyhead[R]{\sffamily\color{HeadingBlue}IB Geography SL}
\fancyfoot[C]{\sffamily\color{HeadingBlue}\thepage\ of \pageref{LastPage}}
\renewcommand{\headrulewidth}{0.4pt}

\newcommand{\term}[1]{\textcolor{HeadingBlue}{\textbf{#1}}}
\newcommand{\exam}[1]{\textcolor{AccentTeal}{\textbf{#1}}}
\newcommand{\todobox}{\fbox{\rule{0pt}{0.8em}\rule{0.8em}{0pt}}}
\newcommand{\smallline}{\rule{0.96\linewidth}{0.4pt}}

\begin{document}

\begin{center}
{\sffamily\bfseries\color{HeadingBlue}\LARGE April 30 Food And Health Essay Planning Workbook}\\[0.25em]
{\sffamily\color{AccentTeal}\large Paper 1 Extended Answer Planning + Underlined Evaluation}
\end{center}

\noindent\colorbox{SoftBlue}{%
\begin{minipage}{0.965\textwidth}
\sffamily
\textbf{\textcolor{HeadingBlue}{Use with the April 30 study pack.}}
The plan requires \term{25 minutes} to plan one \term{Paper 1 extended answer}
from a \term{choice of two prompts}. The special target today is to
\term{underline where evaluation happens}; every main paragraph should include
an underlined evaluative phrase, not only a final conclusion.
\end{minipage}
}

\section{Session Checklist}

\begin{tabularx}{\textwidth}{L{0.12\textwidth}Y L{0.18\textwidth}}
\toprule
\textbf{Done} & \textbf{Task} & \textbf{Time} \\
\midrule
\todobox & Closed-book recall: stakeholders, food security, health inequality,
sustainability & 3 min \\
\todobox & Choose one prompt and unpack command term & 3 min \\
\todobox & Build thesis and criteria & 4 min \\
\todobox & Plan three main paragraphs with evidence and evaluation & 8 min \\
\todobox & Write conclusion in note form with clear judgment & 4 min \\
\todobox & Underline evaluation and fix weakest sentence & 3 min \\
\bottomrule
\end{tabularx}

\section{Minute 0--3: Closed-Book Recall}

Do this before looking back at the study pack.

\begin{tabularx}{\textwidth}{L{0.26\textwidth}Y}
\toprule
\textbf{Recall prompt} & \textbf{My notes} \\
\midrule
Five stakeholders & \smallline \\[0.85em]
Four food-security dimensions & \smallline \\[0.85em]
Three health inequality factors & \smallline \\[0.85em]
One Food and health case & \smallline \\[0.85em]
One sustainability trade-off & \smallline \\[0.85em]
\bottomrule
\end{tabularx}

\section{Minute 3--6: Choose One Prompt}

\noindent\colorbox{SoftTint}{%
\begin{minipage}{0.965\textwidth}
\sffamily
\textbf{\textcolor{HeadingBlue}{Prompt A.}}
\exam{Evaluate} the role of stakeholders in improving food security and health
in one or more places.
\end{minipage}
}

\vspace{0.5em}

\noindent\colorbox{SoftBlue}{%
\begin{minipage}{0.965\textwidth}
\sffamily
\textbf{\textcolor{HeadingBlue}{Prompt B.}}
\exam{To what extent} can sustainable food systems reduce health inequalities?
\end{minipage}
}

\subsection{Question Unpacking}

\begin{tabularx}{\textwidth}{L{0.24\textwidth}Y}
\toprule
\textbf{Decision} & \textbf{Write it here} \\
\midrule
Prompt chosen & A \quad B \\
Command term meaning & \smallline \\[0.8em]
Topic limit & \smallline \\[0.8em]
What I must avoid & \smallline \\[0.8em]
\bottomrule
\end{tabularx}

\section{Minute 6--10: Thesis And Criteria}

\subsection{Criteria For Judgment}

Choose two or three criteria. These are what you will use to judge success,
role, or extent.

\begin{tabularx}{\textwidth}{L{0.1\textwidth}Y Y}
\toprule
\textbf{Use} & \textbf{Criterion} & \textbf{What would count as success?} \\
\midrule
\todobox & Food access & Vulnerable groups can obtain enough safe and
nutritious food \\
\todobox & Health equality & Outcomes improve for poorer, rural, informal, or
otherwise exposed groups \\
\todobox & Diet quality & The response improves nutrition, not only calories or
product availability \\
\todobox & Sustainability & The system remains affordable, lower-waste,
resilient, and resource-conscious over time \\
\todobox & Stakeholder power & Benefits are not controlled by only the most
powerful actor \\
\bottomrule
\end{tabularx}

\subsection{Thesis Builder}

Complete the sentence before planning paragraphs.

\noindent\colorbox{SoftGreen}{%
\begin{minipage}{0.965\textwidth}
\sffamily
Although \rule{0.28\textwidth}{0.4pt}, the strongest judgment is that
\rule{0.48\textwidth}{0.4pt} because \rule{0.33\textwidth}{0.4pt}.
\end{minipage}
}

\section{Minute 10--18: Main Paragraph Plan}

\subsection{Paragraph 1}

\begin{tabularx}{\textwidth}{L{0.23\textwidth}Y}
\toprule
\textbf{Need} & \textbf{Plan notes} \\
\midrule
Point & \smallline \\[0.8em]
Evidence or case & \smallline \\[0.8em]
Explanation chain & \smallline \\[0.8em]
\uline{Evaluation} & \smallline \\[0.8em]
\bottomrule
\end{tabularx}

\subsection{Paragraph 2}

\begin{tabularx}{\textwidth}{L{0.23\textwidth}Y}
\toprule
\textbf{Need} & \textbf{Plan notes} \\
\midrule
Point & \smallline \\[0.8em]
Evidence or case & \smallline \\[0.8em]
Explanation chain & \smallline \\[0.8em]
\uline{Evaluation} & \smallline \\[0.8em]
\bottomrule
\end{tabularx}

\subsection{Paragraph 3}

\begin{tabularx}{\textwidth}{L{0.23\textwidth}Y}
\toprule
\textbf{Need} & \textbf{Plan notes} \\
\midrule
Point & \smallline \\[0.8em]
Evidence or case & \smallline \\[0.8em]
Explanation chain & \smallline \\[0.8em]
\uline{Evaluation} & \smallline \\[0.8em]
\bottomrule
\end{tabularx}

\section{Minute 18--22: Conclusion In Note Form}

Write the conclusion as notes, not full prose. It must answer the exact command
term.

\begin{tabularx}{\textwidth}{L{0.25\textwidth}Y}
\toprule
\textbf{Conclusion need} & \textbf{My notes} \\
\midrule
Overall judgment & \smallline \\[0.9em]
Reason 1 & \smallline \\[0.9em]
Reason 2 & \smallline \\[0.9em]
Limitation or condition & \smallline \\[0.9em]
\bottomrule
\end{tabularx}

\section{Minute 22--25: Evaluation Underline Check}

\rowcolors{2}{TableStripe}{white}
\begin{tabularx}{\textwidth}{L{0.1\textwidth}Y Y}
\toprule
\textbf{Done} & \textbf{Check} & \textbf{Fix if missing} \\
\midrule
\todobox & Thesis makes a judgment, not just a topic statement & Add
"partly", "mainly", "to a large extent", or "limited because" \\
\todobox & Each paragraph has named evidence or a case & Add Brazil and Nestle
or another class-specific Food and health case \\
\todobox & Food security is more precise than "more food" & Name availability,
access, utilization, or stability \\
\todobox & Health inequality is explained through a mechanism & Add income,
gender, age, location, sanitation, healthcare, or retail access \\
\todobox & Evaluation is underlined inside the plan & Add and underline
"however", "more successful for", "less effective where", or "this is limited
because" \\
\bottomrule
\end{tabularx}
\rowcolors{2}{}{}

\subsection{Fix One Evaluative Sentence}

\begin{tabularx}{\textwidth}{Y}
\toprule
\textbf{Weak descriptive sentence} \\
\midrule
\rule{0pt}{1.25cm} \\
\bottomrule
\end{tabularx}

\vspace{0.65em}

\begin{tabularx}{\textwidth}{Y}
\toprule
\textbf{Improved sentence with evaluation underlined} \\
\midrule
\rule{0pt}{1.25cm} \\
\bottomrule
\end{tabularx}

\newpage

\section{Support Sheet}

\subsection{Essay Structure Options}

\rowcolors{2}{TableStripe}{white}
\begin{tabularx}{\textwidth}{L{0.2\textwidth}Y Y}
\toprule
\textbf{Structure} & \textbf{Best for} & \textbf{Paragraph sequence} \\
\midrule
Stakeholder based & Prompt A & Government; TNC or retailer; farmers,
households, NGOs, or health services \\
Food-security based & Either prompt & Availability; access; utilization or
stability; then evaluate which dimension matters most \\
Inequality based & Prompt B & Income and access; health and sanitation; power,
place, and sustainability \\
Sustainability based & Prompt B & Food waste or storage; climate-smart farming
or diets; nexus trade-offs and equity \\
\bottomrule
\end{tabularx}
\rowcolors{2}{}{}

\subsection{Evidence Bank}

\rowcolors{2}{TableStripe}{white}
\begin{tabularx}{\textwidth}{L{0.2\textwidth}Y Y}
\toprule
\textbf{Example or idea} & \textbf{Use it for} & \textbf{Evaluation angle} \\
\midrule
Brazil and Nestle & TNCs, processed food, marketing, distribution, jobs,
nutrition transition, stakeholder conflict & TNCs can improve access and
income opportunities, but may increase processed-food dependence and market
power \\
Government food or health policy & Subsidies, school meals, food safety,
sanitation, healthcare, disease prevention & Can target vulnerable groups, but
is limited by cost, enforcement, targeting, and governance \\
Food-security dimensions & Availability, access, utilization, stability &
More food production does not automatically solve price, diet quality,
sanitation, or shock vulnerability \\
Sustainability futures & Food waste, climate-smart farming, urban agriculture,
healthier diets, water-energy-food nexus & Strategies can reduce long-term
pressure, but benefits depend on cost, scale, culture, and who can participate \\
Disease anchor to verify & Cholera in Haiti, dengue in Singapore, or your
class disease case & Use for diffusion, water or vector pathways, vulnerable
groups, health response, and limits \\
SDG frame & SDG 2, SDG 3, and SDG 6 & Use only if it helps connect hunger,
health, water, sanitation, and uneven progress \\
\bottomrule
\end{tabularx}
\rowcolors{2}{}{}

\subsection{May 1 Case-Bank Starter Choices}

These are candidate anchors to check against class notes, not claims that you
have already studied them.

\rowcolors{2}{TableStripe}{white}
\begin{tabularx}{\textwidth}{L{0.2\textwidth}Y Y}
\toprule
\textbf{Gap} & \textbf{Possible anchor} & \textbf{Use it for} \\
\midrule
Second Food and health case & Cholera in Haiti, dengue in Singapore, a
government school-meal or sanitation policy, or a food-security crisis from
class & Compare public, community, or disease response with the Brazil/Nestle
corporate stakeholder example \\
Climate impact/adaptation case & Bangladesh delta, Sahel drought, European 2003
heatwave, or a class-specific climate case & Turn the archetype into named
evidence: place, hazard, exposed group, response, and limitation \\
Climate response vocabulary & Nature-based solutions / ecosystem-based
adaptation, Paris Agreement / NDCs & Use for response evaluation if the prompt
or stimulus points to ecosystems or international cooperation \\
\bottomrule
\end{tabularx}
\rowcolors{2}{}{}

\subsection{Possible Plan After Timed Attempt}

Use this only after you have made your own plan.

\noindent\colorbox{SoftTint}{%
\begin{minipage}{0.965\textwidth}
\sffamily
\textbf{\textcolor{HeadingBlue}{Prompt A possible judgment.}}
Stakeholders can improve food security and health, but their success is
uneven because governments, TNCs, households, and NGOs control different parts
of the system and do not always share the same priorities.
\end{minipage}
}

\begin{itemize}
  \item Paragraph 1: Governments can improve food safety, sanitation,
  healthcare, school meals, and subsidies; \uline{however, success depends on
  targeting, funding, and governance}.
  \item Paragraph 2: TNCs such as Nestle in Brazil can extend distribution and
  create income opportunities; \uline{however, access to products is not the
  same as improved diet quality or health equality}.
  \item Paragraph 3: NGOs, farmers, and households can support local responses;
  \uline{however, their impact is limited if poverty, prices, or infrastructure
  problems remain unchanged}.
\end{itemize}

\noindent\colorbox{SoftBlue}{%
\begin{minipage}{0.965\textwidth}
\sffamily
\textbf{\textcolor{HeadingBlue}{Prompt B possible judgment.}}
Sustainable food systems can reduce health inequalities to a significant
extent when they improve access, diet quality, and resilience, but they cannot
fully reduce inequality unless poverty, sanitation, healthcare, and stakeholder
power are also addressed.
\end{minipage}
}

\begin{itemize}
  \item Paragraph 1: Food-waste reduction, storage, and better distribution can
  improve availability and stability; \uline{however, this does not guarantee
  affordable household access}.
  \item Paragraph 2: Healthier retail environments, school meals, and public
  health education can improve utilization; \uline{however, choices are still
  shaped by income, time, advertising, and culture}.
  \item Paragraph 3: Climate-smart farming and lower-impact diets can support
  long-term sustainability; \uline{however, poorer farmers and consumers may
  be excluded if costs are high}.
\end{itemize}

\section{Error Log}

\begin{tabularx}{\textwidth}{L{0.24\textwidth}Y Y}
\toprule
\textbf{Weakness} & \textbf{What happened} & \textbf{Fix before May 1} \\
\midrule
Command term & & \\
Case evidence & & \\
Food-security precision & & \\
Health inequality mechanism & & \\
Underlined evaluation & & \\
\bottomrule
\end{tabularx}

\section{Carry Forward}

\noindent\colorbox{SoftGreen}{%
\begin{minipage}{0.965\textwidth}
\sffamily\raggedright
\textbf{\textcolor{HeadingBlue}{For Friday 2026-05-01.}}
The Food and health case detail I must add to the case-study bank is:
\rule{0.30\textwidth}{0.4pt}

The Climate case detail I must add to the case-study bank is:
\rule{0.35\textwidth}{0.4pt}

I have checked whether these are actual class cases before adding them:
\rule{0.28\textwidth}{0.4pt}
\end{minipage}
}

\end{document}
